% Volume VII: Formal Synthesis
% Part XIV. Dynamics, Geometry, and Holonomy
% Chapter 81: Constraint Relaxation
% Spine: proof-spines/ch081-spine.md

\chapter{Constraint Relaxation}
\label{ch:ch081}

This chapter studies how crisis loosens the regularizers that ordinarily discipline institutional revision. Let \(\lambda(t)\) denote procedural regularization. In emergency conditions, \(\lambda(t)\) is reduced, permitting larger updates from limited evidence and relaxing the ordinary barriers that would otherwise constrain motion across belief space. Constraint relaxation can therefore reveal neglected possibilities, but it can also authorize exactly the sort of overreach the earlier chapters were built to resist.

The formal issue is when temporary relaxation becomes durable feature drift at the scale of the state. Lowering \(\lambda(t)\) may be justified when existing procedures are too rigid to register danger or correction, yet the same reduction enlarges the blast radius of weak evidence and weakens the protections embodied in bounded reach, update budgets, and the freedom region. This chapter marks that boundary: disciplined distrust sometimes requires slack, but only on terms that remain reversible and answerable.
